\labsection{5}{Synchronization and interprocess communication}{sec:lab05}
\begin{labproject}{Trace bounded-buffer state and compare a pipe with shared memory}
\textbf{Suggested time:} 100--120 minutes \hfill
\textbf{Level:} intermediate\par
\textbf{Starter files:} \texttt{producer\_consumer\_sim.c}, \texttt{pipe\_demo.c}, \texttt{shm\_writer.c}, \texttt{shm\_reader.c}.\par
\textbf{Evidence to keep:} invariant trace, IPC comparison table, and shared-memory lifecycle evidence.

\begin{labmode}
Two questions stay separate throughout this lab: \emph{how do processes exchange data?} and \emph{how do concurrent activities coordinate safe access to shared state?}
\end{labmode}

\medskip\noindent\textbf{Learning outcomes.}
\begin{enumerate}
  \item use a simple invariant to reason about producer-consumer state;
  \item explain the communication model of a pipe and of shared memory;
  \item identify create/attach/use/detach/remove stages for System V shared memory.
\end{enumerate}

\medskip\noindent\textbf{Part A --- Producer-consumer as a state problem.}
Compile and run \texttt{producer\_consumer\_sim.c}. For each action record \texttt{empty}, \texttt{full}, and whether the action was allowed. With capacity 3, check after every successful operation that
\[
\texttt{empty} + \texttt{full} = 3.
\]
Explain why producing into a full buffer and consuming from an empty buffer must wait or fail rather than blindly changing the counters.

\begin{pitfall}
This first program is a sequential model of semaphore state. It helps you understand the rules, but the integers are not themselves a real thread-safe semaphore implementation.
\end{pitfall}

\medskip\noindent\textbf{Part B --- A pipe moves bytes through the kernel.}
\begin{lstlisting}[style=osbash]
gcc -Wall -Wextra pipe_demo.c -o pipe_demo
./pipe_demo
\end{lstlisting}
Find the call that creates the pipe, the parent's write end, and the child's read end. Explain why the unused ends are closed in each process.

\medskip\noindent\textbf{Part C --- Shared memory exposes the same bytes to two processes.}
\begin{lstlisting}[style=osbash]
gcc -Wall -Wextra shm_writer.c -o shm_writer
gcc -Wall -Wextra shm_reader.c -o shm_reader
ipcs -m
./shm_writer
ipcs -m
./shm_reader
ipcs -m
\end{lstlisting}
Record the shared-memory segment ID and compare the three \texttt{ipcs -m} snapshots. Map each program operation to the lifecycle: create/find $\rightarrow$ attach $\rightarrow$ access $\rightarrow$ detach $\rightarrow$ remove.

\begin{workedexample}
A shared-memory segment is a kernel-managed IPC object. Each process attaches that object into its own address space. Detaching one process does not automatically mean the object has been removed from the system.
\end{workedexample}

\medskip\noindent\textbf{Part D --- Compare the mechanisms.}
Complete a small comparison with rows for data movement, lifetime, setup complexity, and need for synchronization. A pipe provides an ordered byte stream; shared memory gives processes access to a common region and therefore usually needs an additional coordination mechanism when both sides can write.

\medskip\noindent\textbf{Part E --- Failure case.}
Run the reader before the writer and record the error. Then run the normal writer/reader sequence again. A useful systems program should make failure visible instead of silently continuing with invalid state.

\begin{verification}
Submit the producer-consumer trace, one paragraph comparing pipe and shared memory, and terminal evidence showing a shared segment appear and later disappear.
\end{verification}

\begin{troubleshoot}
If a segment remains after a crash, inspect only your own IPC objects with \texttt{ipcs -m}. Remove a known stale segment with \texttt{ipcrm -m ID}. On a shared university machine, never remove another user's IPC object.
\end{troubleshoot}
\end{labproject}
